\documentclass{article}

\usepackage{amsmath,amssymb}
\usepackage{xurl}
\usepackage[hidelinks]{hyperref}

\input{command}

\title{Wolf Chapter~5: Abstract Multiplicative-Domain Hypotheses}
\author{}
\date{2026-07-18}

\begin{document}
\maketitle
\gapnote{clarification}{resolved}

\section{Source statement}

Let $E\colon\MD\to\MD$ be a complex-linear map. Wolf assumes the one-variable
Schwarz inequality
\begin{align}
    E(A^\dagger A)-E(A^\dagger)E(A)&\succeq 0
    \qquad\text{for every $A\in\MD$}.
    \label{eq:wolf_ch5_abstract_multiplicative_domains_schwarz_inequality}
\end{align}
This is Equation~(5.2) of Ref.~\cite{Wolf2012Quantum}.\footnote{Local source:
\path{Notes/WolfNoteTexSource/ch05_schwarz_inequalities.tex},
lines~156--164.}

Wolf defines the right and left multiplicative domains by
\begin{align}
    \mathcal A_R
        &=\left\{A\in\MD:
            E(XA)=E(X)E(A)
            \ \text{for every $X\in\MD$}\right\},
    \label{eq:wolf_ch5_abstract_multiplicative_domains_right_domain}\\
    \mathcal A_L
        &=\left\{A\in\MD:
            E(AX)=E(A)E(X)
            \ \text{for every $X\in\MD$}\right\}.
    \label{eq:wolf_ch5_abstract_multiplicative_domains_left_domain}
\end{align}
These are Equations~(5.11)--(5.12) of
Ref.~\cite{Wolf2012Quantum}.\footnote{Local source:
\path{Notes/WolfNoteTexSource/ch05_schwarz_inequalities.tex},
lines~383--410.}

The corresponding characterizations are
\begin{align}
    A\in\mathcal A_R
        &\quad\Longleftrightarrow\quad
        E(A^\dagger A)=E(A^\dagger)E(A),
    \label{eq:wolf_ch5_abstract_multiplicative_domains_right_characterization}\\
    A\in\mathcal A_L
        &\quad\Longleftrightarrow\quad
        E(AA^\dagger)=E(A)E(A^\dagger).
    \label{eq:wolf_ch5_abstract_multiplicative_domains_left_characterization}
\end{align}
These are Equations~(5.13)--(5.14) of
Ref.~\cite{Wolf2012Quantum}. The theorem assumes complex linearity and
Eq.~\ref{eq:wolf_ch5_abstract_multiplicative_domains_schwarz_inequality}; it
does not assume a Kraus representation, complete positivity, or
$2$-positivity.

\section{Earlier specialization}

The earlier formal characterizations assume a unital Kraus family.\footnote{
Formal declarations:
\leanid{KadisonSchwarz.mem_rightMultiplicativeDomain_iff} and
\leanid{KadisonSchwarz.mem_leftMultiplicativeDomain_iff} in
\doclink{QICLean/Channel/Schwarz/MultiplicativeDomainFull}.}
These declarations are mathematically correct specializations, but their Kraus
and unitality hypotheses do not occur in Wolf's abstract theorem
\cite[Chapter~5, Equations~(5.11)--(5.14)]{Wolf2012Quantum}.

The earlier declarations also use a different naming convention. Their
subscript names the side on which the varying factor is appended. Wolf's
subscript names the side occupied by the fixed element. Consequently, the two
conventions interchange the labels ``left'' and ``right'', although they
describe the same two sets. Their intersections therefore agree:
\begin{align}
    \mathcal A_L\cap\mathcal A_R
        &=\mathcal A_R\cap\mathcal A_L.
    \label{eq:wolf_ch5_abstract_multiplicative_domains_full_domain_agreement}
\end{align}

The source-faithful formalization now assumes only complex linearity and the
Schwarz inequality in
Eq.~\ref{eq:wolf_ch5_abstract_multiplicative_domains_schwarz_inequality}.\footnote{
Formal declarations:
\leanid{IsSchwarzMap},
\leanid{SchwarzMap.rightMultiplicativeDomain},
\leanid{SchwarzMap.leftMultiplicativeDomain},
\leanid{SchwarzMap.mem_rightMultiplicativeDomain_iff}, and
\leanid{SchwarzMap.mem_leftMultiplicativeDomain_iff} in
\doclink{QICLean/Channel/Schwarz/AbstractMultiplicativeDomain}.}
The Kraus declarations remain available as specialized results. They are no
longer used as the blueprint witness for the unrestricted source theorem.

\section{Remaining Chapter~5 boundary}

Wolf's preceding two-variable operator Schwarz inequality and its equality
criterion involve the inverse of $E(B^\dagger B)$ on its range
\cite[Chapter~5, Theorem~5.3]{Wolf2012Quantum}. They are logically separate
from the one-variable result. In particular, the multiplicative-domain theorem
does not assume the two-variable theorem. Its polarization argument applies
Eq.~\ref{eq:wolf_ch5_abstract_multiplicative_domains_schwarz_inequality}
directly to
\begin{align}
    A+tY,
    \qquad
    A,Y&\in\MD,
    \quad
    t\in\C.
    \label{eq:wolf_ch5_abstract_multiplicative_domains_polarization_input}
\end{align}

The two-variable inverse-on-range inequality is treated separately in
\path{docs/paper-gaps/wolf_ch5_two_variable_unconditional.tex}. Its equality
criterion remains separate work and is not needed for the one-variable
multiplicative-domain theorem.

\section{Verdict}

The interchange of the one-sided domain names is a notational clarification.
The two conventions assign opposite labels to the same pair of sets, and their
full multiplicative domains coincide. The source-faithful one-variable theorem
is complete under exactly the abstract hypotheses used by Wolf
\cite[Chapter~5, Equations~(5.2) and (5.11)--(5.14)]{Wolf2012Quantum}. The
two-variable inequality and its equality criterion are separate from this
resolved clarification.

\bibliographystyle{alpha}
\bibliography{references}

\end{document}
